\section{Further results}
\label{app:results}
\begin{figure}[tbp]
\centering\includegraphics[width=\linewidth]{../figures/fig14_caps.pdf}
\caption{Which targets can become strictly first depends on the cap as well as on the rule. Fraction of the mean-non-leading targets of each fixed frame that can be made first, at the smallest number of tasks touched, when each task can rise by at most 5, 10 or 20 points, or by the frame's own task interquartile range (IQR; effective cap never below the target's score), at most 100. The full OLLB frame has the five score rules; Borda, plurality and the Condorcet requirement (dashed, a sufficient condition) are on HELM and the top-500 frame only. Under +10 the minimum rule is the most reachable score rule on full OLLB (446/4,495 vs 113 for the mean); under the IQR cap the sign reverses (283 vs 582). On HELM the minimum is the least reachable score rule at +10 and under IQR. Counts of fixed censuses; the IQR cap is a descriptive data-derived scale, not an achievable training gain. At the ceiling every target is reachable. Source: \texttt{figures/fig14\_caps.provenance.json}.}
\label{fig:capsfig}
\end{figure}
\begin{table}[tbp]
\caption{Targets that can become first (of $N$ mean-non-leading targets) under gains of at most 5, 10 or 20 points per task or the task-IQR cap. Source: \texttt{analysis/results/final\_caps\_feasibility.csv}.}
\label{tab:capscounts}
\centering\scriptsize\setlength{\tabcolsep}{4pt}
\input{../tables/tab_app_caps_counts.tex}
\end{table}
\begin{table}[tbp]
\caption{Frame isolation at cap gain10 (at most $+10$ points per task): fraction of the mean-non-leading targets that can be made strictly first by the mean rule and by the minimum rule, and their paired difference (min $-$ mean). Fixed frames (HELM, full OLLB, OLLB top-500) show the fraction over all their targets; $^\ast$ the bracket is the 2.5--97.5\% target-resampling range (R1000, seed 23) of the paired difference. Matched-frame families keep the 79 models and 78 targets: 100 random 79-model OLLB subsets, 100 random 6-subject HELM subsets (all 79 models), and one fixed OLLB frame of the 79 highest-mean models; entries are median [2.5, 97.5 percentiles] across frames, selected-frame sensitivity ranges, not confidence intervals. $+/-/=$: frames in which the minimum is more, less, equally reachable than the mean. Task corr.: mean pairwise correlation of task scores; gap: mean (best-rival threshold under the minimum rule minus under the mean rule) in score points. Task correlation and gap were supplied only for the matched-frame families; the frames differ in several respects at once, so no single cause is identified. At the ceiling every target is reachable and the difference is 0 in all frames.}
\label{tab:frames}
\centering\scriptsize\setlength{\tabcolsep}{3pt}
\resizebox{\linewidth}{!}{\input{../tables/tab_main_frames.tex}}
\end{table}
\begin{table}[tbp]
\caption{Cap dependence: number of the mean-non-leading targets that can be made strictly first by the mean, minimum and median rules (mean/min/median) when each task can rise by at most 5, 10 or 20 points or by the frame's task interquartile range (IQR). The superscript is the sign of min minus mean ($+$ minimum reaches more targets, $-$ fewer). The full-OLLB advantage of the minimum at $+5$, $+10$, $+20$ reverses under the IQR cap (582 mean vs 283 minimum of 4,495); On top-500 the minimum is ahead at +5, +10 and IQR and behind at +20 (496 vs 499); on HELM it is behind at every cap. Counts of fixed censuses; the IQR cap is a descriptive data-derived scale, not an achievable training gain.}
\label{tab:capsmain}
\centering\scriptsize\setlength{\tabcolsep}{4pt}
\input{../tables/tab_main_caps.tex}
\end{table}
\begin{table}[tbp]
\caption{Matched-frame families, cap ceil: for each rule or requirement, the fraction of the 78 mean-non-leading targets that can be made first, and the median $K$ and $L_1$ over all targets, as the median [2.5, 97.5 percentiles] across the frames of the family (number of frames in brackets after the family). Within a frame the median counts an infeasible target as $+\infty$, so it is $\infty$ when more than half the targets are infeasible; $^\dagger$ median over the feasible targets only (inverted-CDF convention). Condorcet is a requirement (an upper bound). Selected-frame sensitivity ranges, not confidence intervals; the fixed top-79 frame is a single frame. Source: \texttt{analysis/results/final\_frame\_intervals.csv}.}
\label{tab:framegroupsceil}
\centering\scriptsize\setlength{\tabcolsep}{3pt}
\resizebox{\linewidth}{!}{\input{../tables/tab_app_frame_groups_ceil.tex}}
\end{table}
\begin{table}[tbp]
\caption{Matched-frame families, cap gain10: for each rule or requirement, the fraction of the 78 mean-non-leading targets that can be made first, and the median $K$ and $L_1$ over all targets, as the median [2.5, 97.5 percentiles] across the frames of the family (number of frames in brackets after the family). Within a frame the median counts an infeasible target as $+\infty$, so it is $\infty$ when more than half the targets are infeasible; $^\dagger$ median over the feasible targets only (inverted-CDF convention). Condorcet is a requirement (an upper bound). Selected-frame sensitivity ranges, not confidence intervals; the fixed top-79 frame is a single frame. Source: \texttt{analysis/results/final\_frame\_intervals.csv}.}
\label{tab:framegroupsgain10}
\centering\scriptsize\setlength{\tabcolsep}{3pt}
\resizebox{\linewidth}{!}{\input{../tables/tab_app_frame_groups_gain10.tex}}
\end{table}
\begin{figure}[tbp]
\centering\includegraphics[width=\linewidth]{../figures/fig15_units.pdf}
\caption{The mean rule is not the most expensive score rule once costs are measured in task-SD units. Median total investment $L_1$ to make each of the mean-non-leading targets first at the ceiling cap, relative to the mean rule's median, in raw score points (filled) and with each task's rise divided by that task's standard deviation over all models of the frame (open; separately optimised, not a rescaling of the point optimum). Log axis, 1 = mean. On full OLLB the median $L_{1,\mathrm{sd}}$ is 11.1 for the mean rule against 15.4 geometric, 14.4 harmonic, 7.2 median and 6.4 minimum; on HELM it is 49.6 against 50.5, 50.4, 9.9 and 8.2. The point-unit ordering (mean most expensive) and the SD-unit ordering therefore differ. SD units are frame dependent, not training costs. Medians over all targets of fixed censuses; table: values in both units. Source: \texttt{figures/fig15\_units.provenance.json}.}
\label{fig:units}
\end{figure}
\begin{table}[tbp]
\caption{Median costs of the five score rules at the ceiling cap: tasks as a fraction of tasks ($K/n$), points ($L_1$) and largest rise ($L_\infty$) in raw points and in units of each task's standard deviation over all models of the frame (separately optimised, ddof 0). Source: \texttt{analysis/results/final\_costscale\_summary.csv}.}
\label{tab:sdunits}
\centering\scriptsize\setlength{\tabcolsep}{4pt}
\input{../tables/tab_app_costunits.tex}
\end{table}
\begin{figure}[H]
\centering\includegraphics[width=0.8\linewidth]{../figures/fig2_caps.pdf}
\caption{Cap sensitivity (all non-leading targets). (a) feasible fraction; (b, c) medians of $K$ and $L_1$ over feasible targets only, so the target set differs between rules and caps. Caps: $+h$ points per task; e95/e99 empirical gain quantiles (OLLB only, descriptive); fr.: best existing score on each task (+2); 100: ceiling. Missing points: rule not run under that cap. Source: \texttt{figures/fig2\_caps.provenance.json}.}
\label{fig:caps}
\end{figure}
\begin{figure}[H]
\centering\includegraphics[width=0.82\linewidth]{../figures/fig3_paired_ceil.pdf}
\caption{Same-target comparison with the mean rule under the ceiling cap (non-leading targets): share of targets for which the rule is cheaper, equal or costlier than the mean among targets feasible under both. The mean rule is the reference and has no row. Differences mix each rule's gap with its response to increases (the original leader differs), so this is not a gap-matched comparison. Source: \texttt{figures/fig3\_paired\_ceil.provenance.json}.}
\label{fig:paired}
\end{figure}
\begin{figure}[H]
\centering\includegraphics[width=0.82\linewidth]{../figures/fig3_paired_gain10.pdf}
\caption{As Figure~\ref{fig:paired} under the 10-point cap; hatched: feasible under only one of the two rules, white: under neither. Source: \texttt{figures/fig3\_paired\_gain10.provenance.json}.}
\label{fig:paired10}
\end{figure}
\begin{figure}[H]
\centering\includegraphics[width=\linewidth]{../figures/fig1b_feasible_by_rank.pdf}
\caption{Fraction of targets that can become first at all under the 10-point cap, by rank under the mean; Wilson 95\% bands (fixed census). Source: \texttt{figures/fig1b\_feasible\_by\_rank.provenance.json}.}
\label{fig:feasrank}
\end{figure}
\begin{figure}[H]
\centering\includegraphics[width=\linewidth]{../figures/fig6_controls.pdf}
\caption{Targeting controls on HELM (20 targets of mean rank 2--21): optimal $K$ versus random task order (200 orders), and optimal $L_1$ versus the uniform clipped improvement at the optimal $L_\infty$. Source: \texttt{figures/fig6\_controls.provenance.json}.}
\label{fig:controls}
\end{figure}
\begin{figure}[H]
\centering\includegraphics[width=\linewidth]{../figures/fig6c_top500_controls.pdf}
\caption{As Figure~\ref{fig:controls} on the Open LLM Leaderboard top-500 frame (20 targets of mean rank 2--21 within the frame; the frames are not pooled). Source: \texttt{figures/fig6c\_top500\_controls.provenance.json}.}
\label{fig:controls500}
\end{figure}
\begin{figure}[H]
\centering\includegraphics[width=\linewidth]{../figures/fig4_robustness.pdf}
\caption{Task-mixture sensitivity (100 resamplings of the tasks with replacement, cap $+10$, seed 0) on HELM MMLU (57 subjects) and on the Open LLM Leaderboard top-500 frame (6 tasks); ten fixed targets of original mean rank 2--11 within each frame, not pooled. Bands are 2.5--97.5 percentiles, a sensitivity range, not a confidence interval. Source: \texttt{figures/fig4\_robustness.provenance.json}.}
\label{fig:boot}
\end{figure}
\begin{figure}[H]
\centering\includegraphics[width=0.82\linewidth]{../figures/fig7_taskcount.pdf}
\caption{Task-count sensitivity on HELM MMLU: 100 random subsets of $n$ subjects per size, ceiling cap; the target is the model at sorted position 2 or 5 under each rule within each subset, so it differs between rules and subsets. At $n=57$ the subsets differ only by permutation. A sensitivity range, not a confidence interval and not evidence that datasets differ because of their task counts. Source: \texttt{figures/fig7\_taskcount.provenance.json}.}
\label{fig:taskcount}
\end{figure}
\begin{figure}[H]
\centering\includegraphics[width=0.82\linewidth]{../figures/fig8_noise.pdf}
\caption{Sensitivity to assumed Gaussian score noise (independent across models, or correlated 0.8 within a task), cap $+10$, 100 draws; the ten fixed targets of original mean rank 2--11 in each frame. Plug-in standard errors, so a sensitivity range under an assumption, not a measurement-error interval. Source: \texttt{figures/fig8\_noise.provenance.json}.}
\label{fig:noise}
\end{figure}
\begin{figure}[H]
\centering\includegraphics[width=0.82\linewidth]{../figures/fig9_paired_perturb_K.pdf}
\caption{Same-draw comparison with the mean rule under task resampling and score noise, by tasks touched $K$ (ten fixed targets, cap $+10$): share of targets for which the rule needs fewer (filled) or more (open) tasks than the mean on that draw. Ranges over draws, not confidence intervals. The corresponding figure for $L_1$ is Figure~\ref{fig:perturbL1}. Source: \texttt{figures/fig9\_paired\_perturb\_K.provenance.json}.}
\label{fig:perturb}
\end{figure}
\begin{figure}[H]
\centering\includegraphics[width=0.82\linewidth]{../figures/fig9_paired_perturb_L1.pdf}
\caption{As Figure~\ref{fig:perturb} for total points $L_1$. Source: \texttt{figures/fig9\_paired\_perturb\_L1.provenance.json}.}
\label{fig:perturbL1}
\end{figure}
\begin{figure}[H]
\centering\includegraphics[width=\linewidth]{../figures/fig10_paired_bootstrap_ceil.pdf}
\caption{Target-resampling ranges for the same-target comparison with the mean rule under the ceiling cap (1000 resamplings of the targets, seed 23): (a, b) share of targets for which the rule needs fewer tasks (points) than the mean; (c) difference in feasible fraction. Sensitivity ranges of fixed censuses, not population confidence intervals; frames are not pooled. Source: \texttt{figures/fig10\_paired\_bootstrap\_ceil.provenance.json}.}
\label{fig:bootpaired}
\end{figure}
\begin{figure}[H]
\centering\includegraphics[width=\linewidth]{../figures/fig10_paired_bootstrap_gain10.pdf}
\caption{As Figure~\ref{fig:bootpaired} under the 10-point cap. Source: \texttt{figures/fig10\_paired\_bootstrap\_gain10.provenance.json}.}
\label{fig:bootpaired10}
\end{figure}
\begin{figure}[H]
\centering\includegraphics[width=0.82\linewidth]{../figures/fig1c_top500_rank_rules_ceil.pdf}
\caption{Minimum investment versus the target's rank under the mean within the top-500 frame (ceiling cap), per rule including the ordinal rules; the frame is not comparable with the 4,496-model census. Source: \texttt{figures/fig1c\_top500\_rank\_rules\_ceil.provenance.json}.}
\label{fig:top500rank}
\end{figure}
\begin{figure}[H]
\centering\includegraphics[width=0.82\linewidth]{../figures/fig3c_top500_paired_ceil.pdf}
\caption{Same-target comparison with the mean on the top-500 frame, ceiling cap (conventions of Figure~\ref{fig:paired}). The 10-point cap version is Figure~\ref{fig:top500paired10}. Source: \texttt{figures/fig3c\_top500\_paired\_ceil.provenance.json}.}
\label{fig:top500paired}
\end{figure}
\begin{figure}[H]
\centering\includegraphics[width=0.82\linewidth]{../figures/fig3c_top500_paired_gain10.pdf}
\caption{As Figure~\ref{fig:top500paired} under the 10-point cap. Source: \texttt{figures/fig3c\_top500\_paired\_gain10.provenance.json}.}
\label{fig:top500paired10}
\end{figure}
\begin{figure}[H]
\centering\includegraphics[width=0.95\linewidth]{../figures/fig5_theory_check.pdf}
\caption{Solver minimum $L_1$ for the mean against the closed form $n(\tau+\varepsilon-\bar x)$ (Section~\ref{sec:setting}, (i)); the line is the identity. Source: \texttt{figures/fig5\_theory\_check.provenance.json}.}
\label{fig:theory}
\end{figure}
\begin{figure}[t]
\centering\includegraphics[width=\linewidth]{../figures/fig1_rank_rules_ceil.pdf}
\caption{Investment versus the target's rank under the mean (ceiling cap): bin medians with 95\% bootstrap bands over targets in each bin (fixed census, not a population). Ordinal rules only on HELM. Source: \texttt{figures/fig1\_rank\_rules\_ceil.provenance.json}.}
\label{fig:rank}
\end{figure}
\subsection*{Sensitivity: further detail}
\label{app:sens}
Appendix Fig.~\ref{fig:boot} follows ten fixed targets (mean rank 2--11 within each frame) under 100 resamplings of the tasks (cap $+10$). On HELM the number of tasks varies with the draw; the median draw needs fewer tasks under the median rule than under the mean for 8 of the 10 targets; under the minimum rule, whether first place is reachable at all varies strongly across draws (fraction of draws feasible between 0.11 and 1.0 depending on the target, against at least 0.88 for every other rule). On the top-500 Open LLM Leaderboard frame the share of the ten followed targets that both the minimum rule and the mean can win within a draw is lower than on HELM (median over task resamplings 0.4, against 0.7 on HELM). Score-noise perturbations (independent and correlated) and the task-count experiment on HELM ($n=3,\dots,57$ subjects, 100 subsets each) were also run; their percentile ranges are sensitivity ranges under the assumptions of Section~\ref{sec:methods}, not confidence intervals, and are shown in Appendix Figs.~\ref{fig:taskcount}, \ref{fig:noise} and \ref{fig:perturb} (values in analysis/results/summary/). In the task-count experiment the targets are the models at sorted positions 2 and 5 of each random subset, reselected separately for every rule and subset (possible ties are counted in the supplement), so it does not follow one fixed model as $n$ grows and implies no scaling law; for example, the median $K$ of the mean rule for position 2 is 1 at every $n$ from 3 to 57 subjects tested, while its median $L_1$ rises from 1.7 to 7.3 points between $n=3$ and $n=48$ and is 4.0 at $n=57$, where every subset contains all 57 subjects.

\paragraph{Margin and competitor frame.} Two further checks on the 20 followed targets per frame (original mean rank 2--21; Appendix Tables~\ref{tab:eps} and~\ref{tab:frame}; analysis/docs/REVIEW\_FINDINGS.md) show that the costs are conditional on conventions. Changing the winning margin $\varepsilon$ from $0.001$ to $0.1$ changed no feasibility for the mean or Borda at the ceiling or 10-point cap, but changed the task count $K$ of the mean for 2/20 targets at the ceiling and 7/20 at the 10-point cap on HELM (2/20 on the Open LLM Leaderboard), and of Borda for 1/20 on HELM at the 10-point cap; the mean's $L_1$ rises by $n\times0.099$ as the closed form predicts (5.64 points on HELM, 0.59 on the Open LLM Leaderboard). Removing the mean's leader from the rivals (all rules, 10-point cap, original mean rank 2--11) lowered the optimal Borda $K$ for 10/10 HELM targets and the mean's for 5/10, but the minimum rule's for 0/10; under ten random-half rival frames, the share of the ten targets with lower Borda $K$ has median 1.0 (range 0.1--1.0), against 0.5 (0--1.0) for the mean. On this HELM near-top panel, Borda's task count therefore responds more strongly to the rival set than the mean's does. These are near-top panels with ten frames, descriptive ranges and not summaries over all targets.

\subsection*{Headline ranges, median conventions, controls in points and processing choices}
\begin{figure}[H]
\centering\includegraphics[width=0.8\linewidth]{../figures/fig11_policy_controls_L1.pdf}
\caption{As Figure~\ref{fig:controlsK} for total points $L_1$ (the greedy policies lift each touched task to its cap, the last step refined, so these ratios describe those named policies). Source: \texttt{figures/fig11\_policy\_controls\_L1.provenance.json}.}
\label{fig:controlsL1}
\end{figure}
\begin{table}[H]
\caption{Table~\ref{tab:main} with 2.5--97.5 percentile ranges from 1000 resamplings of the targets (seed 23), ceiling cap: median over all non-leading targets (inverted-CDF convention). Descriptive ranges for the fixed censuses, not population confidence intervals. Source: \texttt{analysis/results/final\_headline\_intervals.csv}.}
\label{tab:rangesceil}
\centering\scriptsize\setlength{\tabcolsep}{3pt}
\input{../tables/tab_app_headline_ranges_ceil.tex}
\end{table}
\begin{table}[H]
\caption{As Table~\ref{tab:rangesceil} under the 10-point cap; $\infty$: median over all targets (infeasible targets included); $^\dagger$: median over feasible targets only.}
\label{tab:rangesgain10}
\centering\scriptsize\setlength{\tabcolsep}{3pt}
\input{../tables/tab_app_headline_ranges_gain10.tex}
\end{table}
\begin{table}[H]
\caption{Feasible non-leading targets per score rule, frame and cap, with the number of targets whose optimal cost is zero in parentheses (these enter the ratios in Fig.~\ref{fig:controlsK} as defined in the source file). Source: \texttt{analysis/results/final\_controls\_summary.csv}.}
\label{tab:controlsn}
\centering\scriptsize\setlength{\tabcolsep}{4pt}
\input{../tables/tab_app_controls_n.tex}
\end{table}
\begin{table}[H]
\caption{Entries where the inverted-CDF (lower-middle) and the usual midpoint median differ; the optimisation problem of every target is unchanged. Source: \texttt{analysis/results/final\_headline\_intervals.csv}.}
\label{tab:medconv}
\centering\scriptsize\setlength{\tabcolsep}{3pt}
\input{../tables/tab_app_median_conventions.tex}
\end{table}
\begin{table}[H]
\caption{Share of targets for which each other rule is cheaper than the mean (direction: other rule cheaper), with target-resampling ranges. Source: \texttt{analysis/results/final\_headline\_shares.csv}.}
\label{tab:shares}
\centering\scriptsize\setlength{\tabcolsep}{3pt}
\input{../tables/tab_app_headline_shares.tex}
\end{table}
\begin{figure}[H]
\centering\includegraphics[width=0.9\linewidth]{../figures/fig12_processing.pdf}
\caption{Open LLM Leaderboard processing choices: share of targets that can reach first place at the 10-point cap, by rule. The raw-score cap is a different intervention scale from $+10$ normalised points. Source: \texttt{figures/fig12\_processing.provenance.json}.}
\label{fig:processing}
\end{figure}
\begin{table}[H]
\caption{Reachable targets and ceiling-cap $L_1$ order under six processing choices of the Open LLM Leaderboard frame. Source: \texttt{analysis/results/final\_population\_summary.csv}.}
\label{tab:processing}
\centering\scriptsize\setlength{\tabcolsep}{3pt}
\input{../tables/tab_app_processing.tex}
\end{table}
\subsection*{Further caps and measures}
\begin{table}[H]
\caption{Gain caps of 5 and 20 points per task (same conventions as Table~\ref{tab:main}, including the inverted-CDF median).}
\label{tab:gain520}
\centering\scriptsize\setlength{\tabcolsep}{3pt}
\input{../tables/tab_app_gain5_gain20.tex}
\end{table}
\begin{table}[H]
\caption{Best-existing-score-plus-2 cap and the empirical 95th-percentile gain scale (descriptive only).}
\centering\scriptsize\setlength{\tabcolsep}{3pt}
\input{../tables/tab_app_front2_emp95.tex}
\end{table}
\begin{table}[H]
\caption{Largest single-task rise $L_\infty$, same targets and caps.}
\centering\scriptsize\setlength{\tabcolsep}{3pt}
\input{../tables/tab_app_linf.tex}
\end{table}
\begin{table}[H]
\caption{Targets of original mean rank 2--10 only.}
\centering\scriptsize\setlength{\tabcolsep}{3pt}
\input{../tables/tab_app_rank2_10.tex}
\end{table}
\subsection*{Top-500 frame, margin, rival set and solver audit}
\begin{table}[H]
\caption{Top-500 Open LLM Leaderboard frame (6 tasks), all rules including the ordinal ones; 499 non-leading targets. Conventions of Table~\ref{tab:main}.}
\label{tab:top500}
\centering\scriptsize\setlength{\tabcolsep}{3pt}
\input{../tables/tab_app_top500.tex}
\end{table}
\begin{table}[H]
\caption{Margin sensitivity: $\varepsilon=0.001$ against $0.1$ for the 20 non-leading targets of mean rank 2--21 per frame (not a census).}
\label{tab:eps}
\centering\scriptsize\setlength{\tabcolsep}{3pt}
\input{../tables/tab_app_epsilon.tex}
\end{table}
\begin{table}[H]
\caption{Competitor-frame dependence of $K$ under the 10-point cap for ten fixed targets (mean rank 2--11): the mean's leader removed, and ten random half-frames. A fixed panel, descriptive. The ceiling-cap version is Table~\ref{tab:framec}.}
\label{tab:frame}
\centering\scriptsize\setlength{\tabcolsep}{3pt}
\input{../tables/tab_app_frame_gain10.tex}
\end{table}
\begin{table}[H]
\caption{As Table~\ref{tab:frame} under the ceiling cap.}
\label{tab:framec}
\centering\scriptsize\setlength{\tabcolsep}{3pt}
\input{../tables/tab_app_frame_ceil.tex}
\end{table}
\begin{table}[H]
\caption{Borda solver audit: dynamic program with witness check against the exact MILP fallback, HELM and top-500 instances. Audit counts computed after the main runs, not telemetry of them.}
\label{tab:borda}
\centering\scriptsize\setlength{\tabcolsep}{3pt}
\input{../tables/tab_app_borda_audit.tex}
\end{table}
